\documentclass[
    a4paper,
    DIV=13,
    abstract=true,
    numbers=noenddot
]{scrartcl}

\usepackage{preamble}
\addbibresource{library.bib} % duplicate of preamble, tells VimTeX where entries live

\begin{document}
\title{Gemoetry of Optimal transport}
\author{Liam Llamazares-Elias}
\date{\today}
\maketitle
\section{Introduction}

\section{Classical manifolds}
Clasically a manifold is defined as a space that is locally dipheomorphic to $\R^d$. Smooth functions are those that are locally smooth. 


In $\R^n$ we can think of the tangent space as vectors in $\R^n$. However in manifolds $M$ these arrows do not live in the manifold. As a result we need a differnt definition
\begin{definition}
    The tangent space $T_pM$ is defined as the derivations of the equivalence class of smooth functions identified if they are lcoally equal (germs).
\end{definition}
The definition seems abstract however it has a couple advantages 
\begin{enumerate}
    \item It is defined independantly of coordinates, making it suitable to work with algebraically.
    \item It is isomorphic to $\R^d$. One can still work in terms of coordiantes allowing us to do analysis. A basis is the derivatives with respect to the local coordinates.
    \item It allows us to calculate rates of change (derivatives) of functions in $C^\infty(M)$ and is interpretable in terms of Euclidean space through 
      \begin{aligneq*}
          X_p f= \sum_{j=1}^n v_j \partial_i f = f(\gamma(t)), 
      \end{aligneq*}
      where $\gamma(0)=p$ and $\gamma'(0)=v$. That is we calculate the rate of change of moving $f$ with velocity given by our tangent vecotr $v$.
\end{enumerate}
Of course the greater usefulness is in actually being able to derive theorems and practical from this defintion

Once we deinfe tangent fecotrs at a point (arrows) we can define fields of arrows
\begin{definition}
    A field is a mapping $ p \to X_p \in T_pM$ for all $p \in M$. The field is smooth if it is smooth in coordinate, or equaivalently if $Xf \in C^\infty(M)$ for all $f \in C^\infty(M)$. We write $X \in \Gamma(M)$
\end{definition}
The composition of derivations is not a derivation so we define 
\begin{definition}
  We define the Lie brakcet $[X,Y]= XY-YX \in \Gamma(M)$.
\end{definition}
We also wish to measure the size of our tangent vecotrs
\begin{definition}
    We say $M$ is a Riemannian manifold if there exists a smooth metric $g_p: T_pM \times T_pM \to \R$. That os $g_p$ is PSD for each $p$ and $g(X,Y)$ is smooth.
\end{definition}
Using the metric we can define the gradient. For $\R^n$ we have that
\begin{aligneq*}
  X_p f= \br{ \nabla f ,v}.
\end{aligneq*}
As a result we define 
\begin{definition}
    The gradient of $f$ is the unique element $\nabla f \in T_p M$ such that 
    \begin{aligneq*}
        X_pf = \br{ \nabla f(p), X_p},\quad \forall X \in T_p(M).
    \end{aligneq*}
\end{definition}
The gradient exists by the Riez represeantion therom and it is not just the derivative in each coordinate.

With tangent vectors we can differentiate functions, we extend this definition to differentiate other tangent vectors and $1$-forms.

\begin{definition}
    A connection is $\nabla: \Gamma(M) \times \Gamma(M) \to \Gamma(M)$, written $\nabla(X,Y):=$ such that  
    \begin{enumerate}
      \item $\nabla_{fX} Y= f \nabla_XY$
      \item $\nabla_X (fY)= Xf \nabla_X Y+ f \nabla_XY$.
    \end{enumerate}
\end{definition}
The properties above are reasonable constraints for $\nabla_X$ to actually differntiate $Y$ however there are many such $\nabla$. no a Riemmanina metric there is only one we care about. 

\begin{definition}
    We say that $\nabla$ preserves the metric if 
    \begin{aligneq*}
        X g(Y,Z)= g(\nabla_X Y,Z)+ g(Y, \nabla_XZ).
    \end{aligneq*}
\end{definition}
Yhe trosion measures hte different=ce between moving with $X$ and then $Y$ and the other way around
\begin{definition}
    We say that $\nabla_X$ is torision free if 
    \begin{aligneq*}
      \nabla_XY- \nabla_YX= [X,Y]
    \end{aligneq*}
\end{definition}
\begin{theorem}
    There is a unique conenction that preserves $g$ and is torsion free. We call it the Levi--Civita connection.
\end{theorem}
\begin{definition}
  Given a curve $\gamma$ with $\gamma(0)=p$ and $X(0) \in T_p(M)$  we say that $X(t) \in  T_{\gamma(t)} M$ is the parallel trasnport of $X_0$ along $\gamma$ if 
  \begin{aligneq*}
    \nabla_{\gamma'(t)} X(t)=0
  \end{aligneq*}
\end{definition}
This is many abstract definitions, but how do we justify them?
\begin{theorem}
  Let $P_{t_0, t_1}: T_{\gamma(t_0)}M \to T_{\gamma(t_1)}M$ denote the parallel transport map that takes $V(t_0)$ to $V(t_1)$. Then $\nabla$ preserves the metric iff $P_{t_0,t_1}$ is an isometry.
\end{theorem}
\begin{theorem}
    Symmetry of the HessianLet $f \in C^\infty(M)$. Define the Hessian as the $2$-tensor $\nabla^2 f(X, Y) = X(Yf) - (\nabla_X Y)f$.The connection $\nabla$ is torsion-free if and only if the Hessian is symmetric for all $f$:$$\nabla^2 f(X, Y) = \nabla^2 f(Y, X)$$
\end{theorem}

It is equivalent to definifing thorugh derivative of curves or just $\R^d$.


A Riemannian metric is then a collection of positive definite biliniear form $g_p: T_pM \times T_pM \to \R_+$ such that $g$ is smooth in some sense.



I think this gives some way to define the grad and divergence. Not sure how or why and what spaces they are a part of.Can

\section{Wasserstein Manifold}
We have tha, for quadratic distance, the optimal transport maps take the form 
\begin{aligneq*}
  \set{\nabla \varphi: \varphi \text{ convex }}.
\end{aligneq*}
This suggests that one take the tangent space to be the closure fo this in $L^2$. But why? My interpretation is that one can consider a curve of measure $\mu(t)$. Then, if $\mu(t)$ are geodesically convex one has that 
\begin{aligneq*}
  \mu_{t+\eps}= (I+ \eps T_{\mu_t \to \mu_{t+\eps}})_{\#}\mu_t + \Oo(\eps^2).
\end{aligneq*}
That is, heuristically $\mu_t$ has velocty $\mu'_t =T_{\mu_t \to \mu_{t+\eps}}$. So the velocity is always the grad of a convex function. However this is only for geodesically convex curves. Why consider only these? And this is all heuristic. What properties of normal tangent space for manifolds also hold here?


One shows that AC curves are equal to the solutions of the continuity equation, 
\begin{aligneq*}
    \partial \mu_t =- \nabla \cdot (v_t \mu_t).
\end{aligneq*}
The velocity $v_t$ is not unique so we take the one with minimal $L_2$ norm (this one is unique?). One then defines $\mu'_t:= v_t$ and says that a WGF is a solution to (the CE)
\begin{aligneq*}
  \partial \mu_t = - \nabla_{W_2} \Ff(\mu_t)
\end{aligneq*}
where The WG is defined using the heursitic chain rules
\begin{aligneq*}
  \partial \Ff(\mu_{t})= \br{\nabla_{W_2}\Ff(\mu_t), \mu'_t }
\end{aligneq*}
\liam{It seems to me like there are deeper geometric defintion I am missing so far I don not see the connection to geometry etc. It is defined in terms of equations but I can't see what is happening and how it relates to the tangent of manifolds or the grad of convex functions...}

\section{Generalization to spaces without the continuity equation}
We are interestined in generalig the concep of WGF to spaces where to continuity equation is not well-defined. Concretely we have $X$ as space of grpahs $G=(V,E)$ (some structure on $X$?) and we consider measuers on $X$ and try to build the concept of WGF in this space.

For measures on grpahs this concept exists using the conductivity, however I did not understand there definition of tangent space or see any theorem that showed categorically that this was the right thing to do. Furthermore I di not understand how from the def of tangent space the def of grad flow follwos



\printbibliography
\end{document}
